\section{Introduzione}
Per gran parte della storia della scienza moderna, il riduzionismo ha rappresentato la bussola metodologica predominante: 
l'idea che, per comprendere un sistema, sia sufficiente scomporlo nelle sue parti costituenti e studiarne le leggi fondamentali. Tuttavia,
nel 1972, il Premio Nobel Philip Anderson pubblicò su Science un articolo intitolato "More is Different" \cite{anderson}, 
scardinando questa visione.

Anderson sosteneva che la capacità di ridurre tutto a semplici leggi fondamentali non implica 
la capacità di ricostruire l'universo a partire da esse. Al crescere della complessità e del numero 
di componenti di un sistema, appaiono proprietà totalmente nuove, non deducibili in modo banale dalle 
leggi che governano i singoli elementi. Questo fenomeno è noto come comportamento emergente. I comportamenti emergenti sono proprietà macroscopiche 
di un sistema di cui i singoli componenti non sono portatori. Le proprietà macroscopiche 
non sono deducibili dalle leggi dei singoli, ma emergono dunque dall'interazione di molti elementi.

Un esempio paradigmatico per comprendere il concetto di emergenza è il modello di Ising. Esso descrive un sistema di spin interagenti su 
un reticolo. Ogni spin può assumere due stati, rappresentati da +1 (spin up) e -1 (spin down), e interagisce con i suoi vicini più prossimi.
Il modello di Ising è stato originariamente proposto per studiare il comportamento magnetico dei materiali ferromagnetici, in particolare di studiarne la transizione 
di fase tra uno stato disordinato ad alta temperatura e uno stato ordinato a bassa temperatura. La transizione di fase è un fenomeno emergente: essa non è infatti riconducibile alle proprietà del singolo spin —
il quale non possiede intrinsecamente una temperatura critica o una fase —
ma scaturisce esclusivamente dalla interazione collettive. Sotto la temperatura critica il sistema "sceglie" collettivamente 
una direzione preferenziale di orientazione, trasformando un insieme di interazioni 
locali in un ordine macroscopico manifestato dalla magnetizzazione spontanea.

L’obiettivo di questo lavoro è analizzare la transizione di fase nel modello di Ising attraverso tre prospettive complementari. In primo luogo,
verrà presentata la teoria di campo medio, che fornisce una descrizione analitica approssimata del fenomeno e consente di stimare la temperatura critica.
Successivamente, si studierà il modello tramite simulazioni numeriche Monte Carlo (MC) basate sull’algoritmo di Metropolis,
analizzando grandezze termodinamiche quali energia ($E$), magnetizzazione media ($m$), suscettibilità ($\chi$) e capacità termica ($C_v$)
al variare della temperatura. Inoltrè è stata stimata la temperatura crtica ($T_C$).
Infine, verrà applicata la Principal Component Analysis (PCA) alle configurazioni generate con MC, 
osservando come l'approccio con strumenti di machine learning non supervisionato 
abbia le potenzialità di scoprire transizioni di fase e di riconoscere il parametro d'ordine esclusivamente dalle caratteristiche dei dati.